\section{Experimental Setup}

\subsection{Research Questions}

\begin{enumerate}
    \item \textbf{RQ1 (h-e1)}: Does optimal LoRA rank scale sub-linearly with model size?
    \item \textbf{RQ2 (h-m2)}: Does rank sensitivity increase with model scale?
    \item \textbf{RQ3 (h-c1)}: Is the scaling exponent consistent across task types?
    \item \textbf{RQ4 (h-m1)}: Does attention entropy correlate with model size?
\end{enumerate}

\subsection{Experimental Matrix}

\begin{table}[h]
\centering
\begin{tabular}{lll}
\toprule
\textbf{Dimension} & \textbf{Values} & \textbf{Count} \\
\midrule
Models & Pythia 1B, 2.8B, 6.9B, 12B & 4 \\
Ranks & 4, 8, 16, 32, 64, 128 & 6 \\
Seeds & 42, 1337, 2024 & 3 \\
Tasks & SQuAD-v2, HotpotQA & 2 \\
\midrule
\textbf{Total} & & 144 runs \\
\bottomrule
\end{tabular}
\caption{Experimental matrix dimensions.}
\label{tab:experimental_matrix}
\end{table}

\subsection{Study Scope}

This is a scaling characterization study, not a baseline competition. We measure how optimal rank varies with model size under controlled conditions rather than comparing against alternative methods. The baselines serve as reference points:

\begin{enumerate}
    \item \textbf{Constant rank}: $r=16$ for all model sizes (current practice)
    \item \textbf{Linear scaling}: $r \propto N$ (naive scaling)
    \item \textbf{Full fine-tuning}: All parameters updated (upper bound)
\end{enumerate}
